\documentclass[12pt]{article}
\usepackage{amsmath}
\usepackage{amssymb}
\usepackage[margin=1in]{geometry}

\title{Can a simple model explain economics?}
\author{Ludovico Furlanetto}
\date{June 2026}

\begin{document}
\maketitle

\begin{abstract}
Firm growth shows two robust anomalies. First, volatility falls with size much more
slowly than independence would predict: $\sigma(S)\sim S^{-\beta}$ with
$\beta\approx0.15$--$0.20$ rather than $1/2$. Second, the
growth-rate distribution is sharply peaked and fat-tailed rather than Gaussian.
Standard explanations put the source of these effects inside each firm. We ask
instead whether the same facts can emerge between firms, from a network of
interacting companies in which no heterogeneity is imposed externally. We model the
economy as a Generalized Lotka--Volterra system on a sparse random graph, held at its
cooperative critical point $\mu_c$. This regime is hard to reach numerically: there
the total abundance diverges in finite time and breaks explicit solvers. To get
around this we split the dynamics into a bounded ``shape'' on the probability simplex
and an unbounded ``scale,'' and rescale time so the blow-up turns into smooth
asymptotic growth. The shape can then be integrated stably through the critical
point, and it locates $\mu_c$ cheaply. Using this tool across regular, exponential,
and power-law topologies, we find that the critical dynamics qualitatively reproduce the fat-tailed
growth-rate distribution on the heterogeneous (exponential and power-law) graphs,
though not on the homogeneous one. The size--variance relation, however, is not
reproduced: in this deterministic critical model it is not a stationary law but a
relaxation transient, a V whose minimum is set by the firms' initial sizes rather than
by any intrinsic scale, so the model offers no size--variance exponent to compare with
the data. Static topological criticality thus supplies the heavy tails but not the
slow size--variance decay, which appears to need a genuinely nonstationary mechanism
rather than a static critical one.
\end{abstract}

\end{document}
